\xiaosi

神经电路策略（Neural Circuit Policies, NCP）\cite{lechner2020ncp}的设计灵感来自秀丽隐杆线虫（\textit{C. elegans}）的神经连接组\cite{white1986celegans}。秀丽隐杆线虫仅有302个神经元，却能完成趋化、避障等复杂行为，其关键在于神经元之间高度结构化的稀疏连接模式。NCP将这一思想迁移到人工神经网络中，定义了四层布线架构：

感觉层负责接收外部输入特征，每个输入神经元通过$f_s$条突触向下连接到中间层。中间层对感觉信号进行整合与压缩，每个中间神经元以$f_i$条突触投射到命令层。命令层是决策的核心，不仅接收中间层输入，还具有$r_c$条层内循环连接，使其能够维持短时记忆。运动层作为输出端，每个运动神经元从命令层接收$f_m$条突触输入，生成最终的控制信号。其中$f_s$、$f_i$、$r_c$、$f_m$分别表示感觉扇出、中间扇出、命令循环连接数和运动扇入。

除拓扑结构外，NCP还为每条突触赋予兴奋性（$+1$）或抑制性（$-1$）的极性属性，通过邻接矩阵统一编码。实现中，神经元编号按运动层、命令层、中间层的顺序依次排列。Lechner等人\cite{lechner2020ncp}在车道保持任务上的实验表明，仅19个神经元、253个突触的NCP即可完成任务，而且网络的注意力会自动聚焦到道路边缘等关键区域，无需额外的可解释性方法就能直观理解网络的决策依据。

NCP以LTC单元作为每个神经元的动力学模型，因此继承了LTC的输入自适应特性。与全连接LTC相比，NCP通过稀疏邻接矩阵限制了信息传播路径。这种限制在参数量上是一种压缩，在功能上则赋予了不同层神经元明确的角色分工——感觉层编码、中间层整合、命令层决策、运动层执行——这是全连接网络所不具备的结构性可解释性。
